\section{Further research questions}
\label{questions:section}

The following questions follow directly from proved results or identified
gaps.  They are not presented as solved by the computations in this package.

\subsection{Zero geometry and asymptotic analysis}

\paragraph{1. The first derivative order at which all zeros appear.}
Let $k_0(n)$ be the least $k\geq1$ for which $\gamma_n^{(k)}$ has
$n$ distinct positive zeros.  Our proof shows that $k_0(n)$ exists,
and interlacing shows that saturation persists at every later order.
Determine $k_0(n)$, or obtain an explicit upper bound.
Is $k_0(n)=1$ for every $n$?  This is a concrete question, not a
consequence of our large-$k$ argument.  A useful first project is a
rigorously enclosed zero table for moderate $n$ and all smaller $k$.

\paragraph{2. An effective saturation theorem.}
Turn the proof into a displayed bound $k\geq K(n)$ by bounding the root
separation of $P_n$, the derivatives of $h_n$ near its zeros, and the
gamma concentration remainder.  This would replace a fixed-$n$
existence result by an algorithm with a proved termination bound.
The finite-factor proof of simple roots is qualitative; quantitative
root separation is the main missing input.

\paragraph{3. Joint growth of Laurent index and derivative order.}
The remainder constants currently depend on $n$ without control.
Find a range $n=n(k)\to\infty$ in which the zero count is saturated and
the expansions remain uniform.  This requires a quantitative theory
of the reciprocal-gamma Appell roots, not only fixed-degree
hyperbolicity.  Compare the resulting root distribution with known
asymptotic theory of Jensen and Appell polynomials.

\paragraph{4. Optimal bounds for the first-index root.}
The width-$1/2$ bracket is global and elementary.  Determine the best
constant $C$ for $a_k<e^{H_{k-1}}+C$, or prove sharper two-sided
bounds involving one or two elementary functions of $k$.
The expansion suggests precise limiting offsets, but monotonicity of
those offsets still needs proof.

\paragraph{5. Certified computation of all moving zeros.}
Implement the kernel and the scaled Euler--Maclaurin scheme with ball
arithmetic and explicit tail enclosures.  Combine signs at endpoints
with the proved zero bound to certify completeness.  This would give
an auditable bridge from the analytic theorem to large tables,
including coefficients of high-order zero expansions.

\paragraph{6. Interlacing in the Laurent index.}
The polynomial roots interlace in $n$, and the transformed zeros
eventually inherit the corresponding leading-order ordering.
Determine whether an exact interlacing law holds for
$\gamma_n^{(k)}$ and $\gamma_{n+1}^{(k)}$ at a fixed $k$, and identify
any necessary threshold or counterexample.  Positivity of the kernel
alone does not settle this comparison between different polynomials.

\paragraph{7. Complex zeros and other asymptotic scales.}
Extend the concentration analysis to complex $a$ in sectors.
The real sign-variation bound has no direct complex analogue, so a
complex zero-completeness result requires a separate argument, for
example an argument-principle estimate or a uniform saddle analysis.
Potential transitions near poles of $1/(1-e^{-t})$ should be treated
explicitly rather than inferred from the real theorem.

\subsection{Relations, depth, and exact certificates}

\paragraph{8. Arithmetic injectivity of the formal quotient.}
Specify manageable coefficient fields and background inventories for
which the evaluation map from a weighted distribution quotient is
injective, or has a computable kernel.  The formal rank is now known;
the arithmetic obstruction has consequently been isolated precisely.
Unsuccessful PSLQ searches can suggest test cases but cannot settle
this injectivity problem.

\paragraph{9. Integral presentations and torsion.}
The character proof works over characteristic zero after scalar
extension.  Determine the Smith normal form of the corresponding
integral matrices for positive integer weights, and study denominators
and torsion for negative weights after clearing denominators.
Rational dimension and integral freeness are different questions;
the arbitrary-weight theorem here does not assert the latter.

\paragraph{10. The odd-total-weight Gaussian ladder.}
For $S_{2m}$, compute explicit motivic or coaction classes in a precisely
specified level-four algebra.  The even-weight upper reductions are
proved, but the proposed lower-depth obstruction in the complementary
parity still needs such a calculation.  Start with a small named family
and produce a checkable nonvanishing certificate before asserting a
general alternation theorem.

\paragraph{11. The three-dimensional shuffle quotient at weight six.}
Determine what happens to its three surviving formal directions after
specialization to Gaussian or Eisenstein arguments.  Additional
functional equations, distribution identities, or motivic relations
may lower the dimension.  The exact annihilator witnesses in the
package provide a reproducible starting point, not a proof that the
specialized periods are independent.

\paragraph{12. Complete correction and certificate migration.}
Recompute the modified-polylogarithm complex-embedding tables with the
correct summation endpoint; reconcile Clausen parity conventions;
replace numerical branch guards by exact domain certificates; and
extend the five class-number certificates proved here to the other
fields in subsequent Herglotz investigations. Proving the fitted
Herglotz regulator evaluations remains a separate task even in those
five class-number-one fields. These are well-specified projects
that would make the existing experimental corpus substantially easier
to trust and reuse.

\paragraph{13. Formalization of the elementary core.}
Formalize prime-row factorization, conductor transport, the polynomial
root-preserving lemma, and the Rolle induction for Laplace transforms
in Lean. These reusable elementary results provide manageable foundations
for a later formal treatment of the analytic tower.
